% TODO: 参考文献を以下のように記入．

\begin{thebibliography}{99}

\bibitem{noetel2021}
Noetel, M., et al. (2021).
Video improves learning in higher education: A systematic review.
\textit{Review of Educational Research}, 91(2), 204--236.

\bibitem{liesa2023}
Liesa-Orús, M., et al. (2023).
Digital competence of university teachers: A systematic review.
\textit{Education Sciences}, 13(5), 508.

\bibitem{shoufan2022}
Shoufan, A., \& Mohamed, F. (2022).
YouTube and Education: A Scoping Review.
\textit{IEEE Access}, 10, 125547--125592.

\bibitem{deschaine2015}
Deschaine, M. E., \& Sharma, R. (2015).
The five stages of digital curation: Supporting teaching and learning in online learning environments.
\textit{Journal of Online Learning and Teaching}, 11(3), 398--409.

\bibitem{flintoff2014}
Flintoff, K., Mellow, P., \& Clark, M. (2014).
Digital curation: Opportunities for learning, teaching, research and professional development.
In \textit{Proceedings of ASCILITE 2014}.

\bibitem{hovland1951}
Hovland, C. I., \& Weiss, W. (1951).
The influence of source credibility on communication effectiveness.
\textit{Public Opinion Quarterly}, 15(4), 635--650.

\bibitem{ramlatchan2020}
Ramlatchan, M., \& Watson, G. S. (2020).
Enhancing instructor credibility and immediacy in online multimedia designs.
\textit{Educational Technology Research and Development}, 68, 3239--3264.

\bibitem{bateman2018}
Bateman, J. A., \& Schmidt-Borcherding, F. (2018).
The communicative effectiveness of education videos:
Towards an empirically motivated multimodal account.
\textit{Multimodal Technologies and Interaction}, 2(3), 59.

\end{thebibliography}